% Bu dosya, Bolum 10'u zenginlestirir (Yeni chapter acmaz)
\section{Kontrol Haritalama ve Denetim Hazırlığı}

\subsection{ISO 27001 – NIST CSF – CIS Controls Eşlemesi}
\begin{longtable}{|p{3.5cm}|p{5.5cm}|p{5.5cm}|}
\hline
\rowcolor{tableheadcolor}
\textbf{Kontrol Alanı} & \textbf{ISO 27001 Annex A (ör.)} & \textbf{NIST CSF/CIS (ör.)} \\
\hline
Erişim Yönetimi & A.9 Erişim Kontrolü & PR.AC; CIS 6, 16 \\
\hline
Varlık Yönetimi & A.8 Varlıkların Yönetimi & ID.AM; CIS 1, 2 \\
\hline
Kriptografi & A.10 Kriptografi & PR.DS; CIS 3 \\
\hline
Operasyon Güvenliği & A.12 Operasyonel Güvenlik & PR.IP; CIS 4, 8 \\
\hline
Tedarikçi Güvenliği & A.15 Tedarikçi İlişkileri & ID.SC; CIS 15 \\
\hline
Olay Yönetimi & A.16 Bilgi Güvenliği İhlalleri & DE/RS; CIS 17 \\
\hline
\end{longtable}

\subsection{Olgunluk Düzeyi (CMMI benzeri) Önerisi}
\begin{itemize}
  \item \textbf{Seviye 1 – Ad-hoc:} Reaktif, kişiye bağlı
  \item \textbf{Seviye 2 – Tanımlı:} Politika/prosedür var, düzensiz uygulama
  \item \textbf{Seviye 3 – Uygulanan:} Standartlaştırılmış uygulama ve temel metrikler
  \item \textbf{Seviye 4 – Ölçümlenen:} KPI/KRI tanımlı, sürekli ölçüm
  \item \textbf{Seviye 5 – Optimize:} Sürekli iyileştirme, otomasyon
\end{itemize}

\subsection{Denetim Hazırlık Kontrol Listesi}
\begin{itemize}
  \item Varlık/envanter ve veri sınıflandırma güncel mi?
  \item Politika/prosedürler versiyonlanmış ve duyurulmuş mu?
  \item IAM ve erişim gözden geçirmeleri periyodik mi? (RBAC/ABAC)
  \item Olay müdahale planı ve tatbikat kayıtları mevcut mu?
  \item Üçüncü taraf risk değerlendirmeleri güncel mi?
\end{itemize}

\Not{Kontrol eşlemesini tek bir “kaynak gerçekliği” tablosunda tutun. Denetimlerde çelişen versiyonları engeller.}

